% ---------------------------------------------------------------------------
% Anhang D: Testunterlagen
% Erzeugt aus Testmaterial/_generator (aufgaben.py, werte.py, p4_druck.py,
% BEZUG aus make_docs.py) mit thesis/werkzeuge/anhang_d.py.
% Nicht von Hand aendern, sondern neu erzeugen.
% ---------------------------------------------------------------------------

\chapter{Testunterlagen}
\label{anh:testunterlagen}

Die Unterlagen stehen hier im Wortlaut der gedruckten Fassungen, die in den elf Sitzungen verwendet wurden (Abschnitt~\ref{sec:testdesign}). Jede Aufgabenkarte war eine eigene Seite, auf der allein der Kartentext stand; Startzustand und Erfolgskriterium kannte nur die Testleitung. Zeitbudget und Kennzahlen je Aufgabe nennt Tabelle~\ref{tab:aufgaben}.

\section*{Aufgabenkarten}

\begingroup
\scriptsize
\setlength{\tabcolsep}{3pt}
\renewcommand{\arraystretch}{1.12}
\setlength{\LTcapwidth}{\textwidth}
\begin{longtable}{@{}L{0.6cm}L{6.1cm}L{3.8cm}L{4.2cm}@{}}
	\caption[Aufgabenkarten Block A: ohne Anmeldung]{Aufgabenkarten Block A: ohne Anmeldung, Kartentext wortgleich mit Startzustand und Erfolgskriterium der Testleitung (eigene Darstellung)}
	\label{tab:karten-a}\\
	\toprule
	\textbf{ID} & \textbf{Kartentext} & \textbf{Startzustand} & \textbf{Erfolgskriterium} \\
	\midrule
	\endfirsthead
	\toprule
	\textbf{ID} & \textbf{Kartentext} & \textbf{Startzustand} & \textbf{Erfolgskriterium} \\
	\midrule
	\endhead
	\bottomrule
	\endfoot
	A1 & „Du bist gerade auf der Startseite dieses Portals gelandet, ohne dass dir jemand erklärt hat, was das ist. Schau dich um und sag laut, was du vermutest: Was ist das für eine Seite, wer betreibt sie, und was kann man hier tun?“ & Startseite, abgemeldet, frisches Browserprofil, kein Scrollen vorbereitet. Die Top 10 ist vollständig belegt. & Keine Erfolgswertung und keine Zeitnahme gegen ein Kriterium. Erhoben wird, ob Zweck, Betreiber und mindestens eine Handlungsmöglichkeit ohne Hilfe benannt werden. \\
	A2 & „Eine Bekannte hat dir erzählt, dass an der Hochschule ein Projekt zum Thema Baustoffbörse entstanden ist. Finde dieses Projekt.“ & Startseite, abgemeldet. & Detailseite des Zielprojekts ist geöffnet, Hilfestufe höchstens 2. \\
	A3 & „Jetzt möchtest du dir einen Überblick verschaffen. Zeig mir alle Bachelorarbeiten der Fakultät Elektrotechnik und Informatik aus dem Semester WiSe 2025/26.“ & Projektliste, keine Filter gesetzt, abgemeldet. & Alle drei Einschränkungen gleichzeitig aktiv (Abschluss, Fakultät, Semester), Trefferliste sichtbar. \\
	A4 & „Die Liste ist dir jetzt zu eng geworden. Erweitere sie wieder, ohne von vorn anzufangen. Öffne danach ein Projekt, das dich interessiert, und finde heraus, was dieselbe Person sonst noch veröffentlicht hat.“ & Ergebnis aus A3 unverändert stehen lassen. & Mindestens eine Einschränkung entfernt, ohne Neuladen der Seite und ohne Rückkehr zur Startseite; danach ist die Portfolio-Profilseite der Autorin oder des Autors geöffnet. \\
	A5 & „Du willst dieses Projekt jemandem zeigen, der gerade kein Internet hat. Sichere es dir so, dass du es weitergeben kannst.“ & Das Projekt „Mehrwegpfand in der Bremer Gastronomie“ (aus A2 oder A4) geöffnet stehen lassen, abgemeldet. & Eine PDF-Datei ist erzeugt und gespeichert, ohne Hilfestufe 3. \\
\end{longtable}
\endgroup

\begingroup
\scriptsize
\setlength{\tabcolsep}{3pt}
\renewcommand{\arraystretch}{1.12}
\setlength{\LTcapwidth}{\textwidth}
\begin{longtable}{@{}L{0.6cm}L{6.1cm}L{3.8cm}L{4.2cm}@{}}
	\caption[Aufgabenkarten Block B: Studierende]{Aufgabenkarten Block B: Studierende, Kartentext wortgleich mit Startzustand und Erfolgskriterium der Testleitung (eigene Darstellung)}
	\label{tab:karten-b}\\
	\toprule
	\textbf{ID} & \textbf{Kartentext} & \textbf{Startzustand} & \textbf{Erfolgskriterium} \\
	\midrule
	\endfirsthead
	\toprule
	\textbf{ID} & \textbf{Kartentext} & \textbf{Startzustand} & \textbf{Erfolgskriterium} \\
	\midrule
	\endhead
	\bottomrule
	\endfoot
	B1 & „Hier sind deine Zugangsdaten. Melde dich am Portal an.“ \newline \textit{Darunter je Sitzung: E-Mail und Passwort des Testkontos.} & Abgemeldet, Projektliste geöffnet. Auf der Karte stehen die Anmeldeadresse und das Passwort. & Angemeldet, höchstens zwei Fehlversuche. \\
	B2 & „Du möchtest deine Arbeit aus dem letzten Semester im Portal zeigen. Lege sie hier an. Sie heißt ‚Lastenrad-Verleih im Bremer Westen‘.“ & Angemeldet als Studierende, Startseite. & Ein Projekt im Status Entwurf existiert und trägt den vorgegebenen Titel. \\
	B3 & „Baue die erste Seite deines Projekts: eine Überschrift, einen Absatz Text, ein Bild und eine Aufzählung mit drei Punkten. Was genau dort steht, ist egal.“ & Leerer Editor des in B2 angelegten Projekts; eine Bilddatei liegt sichtbar auf dem Schreibtisch des Testrechners. & Überschrift, Absatz, Bild und Aufzählung sind im Dokument, die Speicherung ist bestätigt. \\
	B4 & „Dein Projekt soll zur Prüfung eingereicht werden. Trage alles ein, was dafür nötig ist.“ & Ergebnis aus B3 unverändert stehen lassen: Projekt mit Inhalt, keine Metadaten gesetzt. & Alle sieben Pflichtangaben gesetzt: Projekttyp, Fakultät, Studiengang, Abschluss, Kategorie, Semester und mindestens eine zugewiesene Redaktion. \\
	B5 & „Gib dein Projekt zur Prüfung ab.“ & Ergebnis aus B4 unverändert stehen lassen: Projekt vollständig, Status Entwurf. & Status Eingereicht. \\
	B6 & „Dir fällt ein, dass ein Absatz fehlt. Sorge dafür, dass du wieder weiterschreiben kannst.“ & Ergebnis aus B5 unverändert stehen lassen: Projekt im Status Eingereicht. & Projekt wieder im Status Entwurf, Editor geöffnet. \\
	B7 & „Zum Schluss: Ergänze in deinem eigenen Profil eine kurze Beschreibung von dir, damit Besucher wissen, wer hinter den Projekten steht.“ & Angemeldet, eigenes Profil ohne Kurzbiografie. & Kurzbiografie gespeichert. \\
\end{longtable}
\endgroup

\begingroup
\scriptsize
\setlength{\tabcolsep}{3pt}
\renewcommand{\arraystretch}{1.12}
\setlength{\LTcapwidth}{\textwidth}
\begin{longtable}{@{}L{0.6cm}L{6.1cm}L{3.8cm}L{4.2cm}@{}}
	\caption[Aufgabenkarten Block C: Redaktion]{Aufgabenkarten Block C: Redaktion, Kartentext wortgleich mit Startzustand und Erfolgskriterium der Testleitung (eigene Darstellung)}
	\label{tab:karten-c}\\
	\toprule
	\textbf{ID} & \textbf{Kartentext} & \textbf{Startzustand} & \textbf{Erfolgskriterium} \\
	\midrule
	\endfirsthead
	\toprule
	\textbf{ID} & \textbf{Kartentext} & \textbf{Startzustand} & \textbf{Erfolgskriterium} \\
	\midrule
	\endhead
	\bottomrule
	\endfoot
	C1 & „Hier sind deine Zugangsdaten. Melde dich am Portal an.“ \newline \textit{Darunter je Sitzung: E-Mail und Passwort des Testkontos.} & Abgemeldet, Projektliste geöffnet. Auf der Karte stehen die Anmeldeadresse und das Passwort. & Angemeldet, höchstens zwei Fehlversuche. \\
	C2 & „Du bist in der Redaktion für einige eingereichte Projekte zuständig. Öffne deine Reviews und verschaffe dir einen Überblick, was gerade auf dich wartet.“ & Angemeldet als Redaktion, Startseite; vier Projekte warten im Arbeitsvorrat. Nach C1 steht „Meine Reviews“ schon offen: die Moderation öffnet die Startseite, bevor sie die Karte reicht. & Die Liste der eigenen Reviews ist geöffnet, mindestens ein wartendes Projekt wird benannt. \\
	C3 & „In deinen Reviews kannst du an jedem Projekt drei Arten von Kommentaren schreiben. Eine Antwort ist eine Rückfrage, die alle lesen, die das Projekt sehen dürfen. Eine Notiz ist intern: nur Redaktion und Administration sehen sie. Eine Änderung fordert die Person, die das Projekt eingereicht hat, auf, etwas zu überarbeiten, und gibt ihr das Projekt dafür zurück. Deine Aufgabe: Öffne das Projekt ‚Zahlungsbereitschaft für regionale Milch in Bremen‘. Dir fällt auf, dass die Quellenangaben fehlen. Sorge dafür, dass der Autor das nachträgt.“ & Seed-Projekt „Zahlungsbereitschaft für regionale Milch in Bremen“ im Status Eingereicht, die Testperson ist ihm als Redaktion zugewiesen. & Status Überarbeitung, Begründung gespeichert und für den Autor sichtbar. \\
	C4 & „Auf dem Weg zur Veröffentlichung durchläuft ein Projekt vier Stände. Wer es schreibt, legt es als Entwurf an und reicht es ein, dann steht es auf Eingereicht. Die Redaktion prüft und genehmigt es, dann steht es auf Genehmigt. Erst wenn die Redaktion es danach veröffentlicht, steht es auf Veröffentlicht und ist für alle sichtbar. Deine Aufgabe: Das Projekt ‚Wiederverwendung von Ziegeln aus einem Abbruch in Walle‘ ist in Ordnung. Veröffentliche es.“ & Zweites Seed-Projekt „Wiederverwendung von Ziegeln aus einem Abbruch in Walle“ im Status Eingereicht, die Testperson ist ihm als Redaktion zugewiesen. & Status Veröffentlicht, ohne Hilfestufe 3. \\
	C5 & „Das Projekt, das du gerade veröffentlicht hast, gefällt dir besonders gut. Nominiere es für die Top 10 auf der Startseite.“ & Das in C4 veröffentlichte Projekt „Wiederverwendung von Ziegeln aus einem Abbruch in Walle“ ist noch in der Prüfansicht geöffnet; es trägt keine offene Nominierung und keinen Top-10-Platz. In der Top 10 sind sechs der zehn Plätze belegt. & Nominierung mit Stufe ist abgegeben. \\
	C6 & „Mit dem veröffentlichten Projekt bist du fertig. Öffne jetzt ein anderes Projekt, ‚Audiowalk zur Werftgeschichte in Gröpelingen‘, und halte daran eine Notiz fest, die nur die Redaktion sehen soll und nicht die Autorin.“ & Drittes Seed-Projekt „Audiowalk zur Werftgeschichte in Gröpelingen“ im Status Eingereicht, die Testperson ist ihm als Redaktion zugewiesen. Geöffnet ist noch das in C4 veröffentlichte und in C5 nominierte Projekt. & Eine Notiz ist angelegt, keine Antwort und keine Änderung. \\
\end{longtable}
\endgroup

\begingroup
\scriptsize
\setlength{\tabcolsep}{3pt}
\renewcommand{\arraystretch}{1.12}
\setlength{\LTcapwidth}{\textwidth}
\begin{longtable}{@{}L{0.6cm}L{6.1cm}L{3.8cm}L{4.2cm}@{}}
	\caption[Aufgabenkarten Block D: Administration]{Aufgabenkarten Block D: Administration, Kartentext wortgleich mit Startzustand und Erfolgskriterium der Testleitung (eigene Darstellung)}
	\label{tab:karten-d}\\
	\toprule
	\textbf{ID} & \textbf{Kartentext} & \textbf{Startzustand} & \textbf{Erfolgskriterium} \\
	\midrule
	\endfirsthead
	\toprule
	\textbf{ID} & \textbf{Kartentext} & \textbf{Startzustand} & \textbf{Erfolgskriterium} \\
	\midrule
	\endhead
	\bottomrule
	\endfoot
	D1 & „Hier sind deine Zugangsdaten. Melde dich am Portal an.“ \newline \textit{Darunter je Sitzung: E-Mail und Passwort des Testkontos.} & Abgemeldet, Projektliste geöffnet. Auf der Karte stehen die Anmeldeadresse und das Passwort. & Angemeldet, höchstens zwei Fehlversuche. \\
	D2 & „Du betreust dieses Portal. Öffne die Administrationsseite und verschaff dir dort einen Überblick, was gerade im Portal los ist.“ & Angemeldet als Administration, Startseite des Portals (nicht der Administrationsseite). Nach D1 steht die Administrationsseite schon offen: die Moderation öffnet die Startseite, bevor sie die Karte reicht. & Die Administrationsseite ist geöffnet, und mindestens zwei konkrete Aussagen über den Zustand des Portals werden genannt. \\
	D3 & „Nele Ostermann soll künftig Projekte prüfen dürfen. Ändere dafür die Rolle des Kontos auf Redakteur:in.“ & Administrationsseite offen; das Zielkonto trägt die Rolle STUDENT und eine studentische HSB-Adresse (@stud.hs-bremen.de). Die Hochstufung ist der Tutor-Fall aus Contract §3.1.1: eine studentische Adresse belegt den Studierendenstatus, nicht die redaktionelle Tätigkeit. & Das Konto trägt die Rolle EDITOR. \\
	D4 & „Die Kategorien ‚Studienprojekt‘ und ‚Studentisches Projekt‘ meinen dasselbe. Führe die beiden zu einer Kategorie zusammen, ohne dass Projekte verloren gehen.“ & Die zwei gleichbedeutenden Kategorien „Studienprojekt“ und „Studentisches Projekt“ mit je mindestens drei Projekten daran. Daneben stehen ähnlich klingende Stammkategorien wie „Semesterprojekt“; deshalb nennt die Karte beide Namen. & Nur noch eine der beiden Kategorien existiert, und alle betroffenen Projekte hängen daran. \\
	D5 & „Die fünf genehmigten Projekte sollen veröffentlicht werden. Erledige das.“ & Projektverwaltung offen, fünf Projekte im Status Genehmigt, über den Statusfilter auffindbar. & Alle fünf Projekte tragen den Status Veröffentlicht. \\
	D6 & „Die Redaktion hat Projekte für die Top 10 der Startseite vorgeschlagen. Nimm aus diesen Vorschlägen Projekte in die Top 10 auf, bis alle zehn Plätze belegt sind.“ & Mindestens zwölf offene Nominierungen, keine Top 10 gesetzt. & Alle zehn Plätze sind belegt; die Verwaltung zeigt „10 von 10 Plätzen belegt“. Einen eigenen Speicherschritt gibt es nicht. \\
	D7 & „Jemand behauptet, die Top 10 der Startseite sei vor deiner Arbeit daran schon einmal geändert worden. Finde heraus, wann das zuletzt war und wer es war.“ & Ergebnis aus D6 unverändert stehen lassen: Top 10 gesetzt, im Änderungsprotokoll stehen mehrere Einträge von verschiedenen Personen. & Das Änderungsprotokoll ist geöffnet; die eigenen Einträge aus D6 werden als solche erkannt, und der jüngste Eintrag davor wird mit Person und Zeitpunkt genannt. \\
\end{longtable}
\endgroup

\section*{Fragebogen zur Gebrauchstauglichkeit (SUS)}

Deutsche Übersetzung über dem englischen Original nach Brooke (1996). Einleitung auf dem Bogen: „Mit ‚dem System‘ ist in allen zehn Aussagen das Portal gemeint, das du gerade benutzt hast. Bitte kreuze in jeder Zeile genau ein Feld an. Es gibt kein richtig und kein falsch.“ Skala: 1 = stimme überhaupt nicht zu, 5 = stimme voll zu.

\begingroup
\footnotesize
\setlength{\tabcolsep}{4pt}
\renewcommand{\arraystretch}{1.15}
\setlength{\LTcapwidth}{\textwidth}
\begin{longtable}{@{}L{0.7cm}L{9.6cm}*{5}{>{\centering\arraybackslash}p{0.62cm}}@{}}
	\caption[SUS-Bogen]{SUS-Bogen in der eingesetzten deutschen Fassung mit englischem Original (nach Brooke, 1996)}
	\label{tab:sus-bogen}\\
	\toprule
	\textbf{Nr.} & \textbf{Aussage} & \textbf{1} & \textbf{2} & \textbf{3} & \textbf{4} & \textbf{5} \\
	\midrule
	\endfirsthead
	\toprule
	\textbf{Nr.} & \textbf{Aussage} & \textbf{1} & \textbf{2} & \textbf{3} & \textbf{4} & \textbf{5} \\
	\midrule
	\endhead
	\bottomrule
	\endfoot
	1 & Ich denke, dass ich das System gerne regelmäßig nutzen würde. \newline {\scriptsize\itshape I think that I would like to use this system frequently} & $\square$ & $\square$ & $\square$ & $\square$ & $\square$ \\
	2 & Ich fand das System unnötig komplex. \newline {\scriptsize\itshape I found the system unnecessarily complex} & $\square$ & $\square$ & $\square$ & $\square$ & $\square$ \\
	3 & Ich fand das System einfach zu benutzen. \newline {\scriptsize\itshape I thought the system was easy to use} & $\square$ & $\square$ & $\square$ & $\square$ & $\square$ \\
	4 & Ich glaube, ich würde die Unterstützung einer technisch versierten Person benötigen, um das System zu benutzen. \newline {\scriptsize\itshape I think that I would need the support of a technical person to use this system} & $\square$ & $\square$ & $\square$ & $\square$ & $\square$ \\
	5 & Ich fand, die verschiedenen Funktionen des Systems waren gut integriert. \newline {\scriptsize\itshape I found the various functions in this system were well integrated} & $\square$ & $\square$ & $\square$ & $\square$ & $\square$ \\
	6 & Ich denke, das System enthielt zu viele Inkonsistenzen. \newline {\scriptsize\itshape I thought there was too much inconsistency in this system} & $\square$ & $\square$ & $\square$ & $\square$ & $\square$ \\
	7 & Ich kann mir vorstellen, dass die meisten Menschen den Umgang mit diesem System sehr schnell lernen. \newline {\scriptsize\itshape I would imagine that most people would learn to use this system very quickly} & $\square$ & $\square$ & $\square$ & $\square$ & $\square$ \\
	8 & Ich fand das System sehr umständlich zu nutzen. \newline {\scriptsize\itshape I found the system very cumbersome to use} & $\square$ & $\square$ & $\square$ & $\square$ & $\square$ \\
	9 & Ich fühlte mich bei der Nutzung des Systems sehr sicher. \newline {\scriptsize\itshape I felt very confident using the system} & $\square$ & $\square$ & $\square$ & $\square$ & $\square$ \\
	10 & Ich musste eine Menge lernen, bevor ich mit dem System zurechtkam. \newline {\scriptsize\itshape I needed to learn a lot of things before I could get going with this system} & $\square$ & $\square$ & $\square$ & $\square$ & $\square$ \\
\end{longtable}
\endgroup

\section*{Fragebogen zur Einschätzung des Portals (TAM)}

Deutsche Übersetzung über dem englischen Original nach Davis (1989); [X] steht im Original für den Namen des Systems. Für BI1 und BI2 stand kein englisches Original auf dem Bogen. Der Bogen lag in vier Fassungen vor, die sich allein im Bezugssatz über den Aussagen unterscheiden:

\begin{itemize}[topsep=0.3em, itemsep=0.2em]
	\item Fassung A (Studentische Projekte suchen und ansehen): „Denk beim Beantworten an das Suchen und Ansehen studentischer Projekte, so wie du es eben getan hast.“
	\item Fassung B (Eigene Projektarbeiten erstellen und einreichen): „Denk beim Beantworten an das Erstellen und Einreichen eigener Projektarbeiten, so wie du es eben getan hast.“
	\item Fassung C (Eingereichte Projekte prüfen und veröffentlichen): „Denk beim Beantworten an das Prüfen und Veröffentlichen eingereichter Projekte, so wie du es eben getan hast.“
	\item Fassung D (Das Portal verwalten): „Denk beim Beantworten an die Verwaltung des Portals, so wie du es eben getan hast.“
\end{itemize}

Einleitung auf dem Bogen: „Vierzehn Aussagen. Kreuze in jeder Zeile genau ein Feld an, je nachdem, wie wahrscheinlich die Aussage für dich ist. Es gibt keine richtigen und keine falschen Antworten.“ Die Felder trugen auf dem Bogen keine Zahlen, sondern Beschriftungen; in Teil 1, Teil 2 und bei BI1 von links: äußerst wahrscheinlich, ziemlich wahrscheinlich, etwas wahrscheinlich, weder noch, etwas unwahrscheinlich, ziemlich unwahrscheinlich, äußerst unwahrscheinlich. Kodiert wurde nach Position, das linke Feld mit 7 und das rechte mit 1, sodass hohe Werte Zustimmung bedeuten; die Spaltenköpfe der Tabelle geben diese Kodierung wieder. Bei BI2 sind die Endpunkte vertauscht gedruckt (links nicht zu erwarten, rechts zu erwarten); der Wert ging deshalb als 8 minus Rohwert ein.

\begingroup
\footnotesize
\setlength{\tabcolsep}{3pt}
\renewcommand{\arraystretch}{1.15}
\setlength{\LTcapwidth}{\textwidth}
\begin{longtable}{@{}L{1.15cm}L{8.45cm}*{7}{>{\centering\arraybackslash}p{0.5cm}}@{}}
	\caption[TAM-Bogen]{TAM-Bogen in der eingesetzten deutschen Fassung mit englischem Original (nach Davis, 1989)}
	\label{tab:tam-bogen}\\
	\toprule
	\textbf{Item} & \textbf{Aussage} & \textbf{7} & \textbf{6} & \textbf{5} & \textbf{4} & \textbf{3} & \textbf{2} & \textbf{1} \\
	\midrule
	\endfirsthead
	\toprule
	\textbf{Item} & \textbf{Aussage} & \textbf{7} & \textbf{6} & \textbf{5} & \textbf{4} & \textbf{3} & \textbf{2} & \textbf{1} \\
	\midrule
	\endhead
	\bottomrule
	\endfoot
	\multicolumn{9}{@{}l}{\textbf{Teil 1: Nutzen für deine Aufgaben}} \\
	PU1 & Die Nutzung des Portals würde es mir ermöglichen, diese Aufgaben schneller zu erledigen. \newline {\scriptsize\itshape Using [X] in my job would enable me to accomplish tasks more quickly.} & $\square$ & $\square$ & $\square$ & $\square$ & $\square$ & $\square$ & $\square$ \\
	PU2 & Die Nutzung des Portals würde meine Leistung bei diesen Aufgaben verbessern. \newline {\scriptsize\itshape Using [X] would improve my job performance.} & $\square$ & $\square$ & $\square$ & $\square$ & $\square$ & $\square$ & $\square$ \\
	PU3 & Die Nutzung des Portals würde meine Produktivität bei diesen Aufgaben erhöhen. \newline {\scriptsize\itshape Using [X] in my job would increase my productivity.} & $\square$ & $\square$ & $\square$ & $\square$ & $\square$ & $\square$ & $\square$ \\
	PU4 & Die Nutzung des Portals würde meine Effektivität bei diesen Aufgaben steigern. \newline {\scriptsize\itshape Using [X] would enhance my effectiveness on the job.} & $\square$ & $\square$ & $\square$ & $\square$ & $\square$ & $\square$ & $\square$ \\
	PU5 & Die Nutzung des Portals würde mir diese Aufgaben erleichtern. \newline {\scriptsize\itshape Using [X] would make it easier to do my job.} & $\square$ & $\square$ & $\square$ & $\square$ & $\square$ & $\square$ & $\square$ \\
	PU6 & Ich würde das Portal für diese Aufgaben nützlich finden. \newline {\scriptsize\itshape I would find [X] useful in my job.} & $\square$ & $\square$ & $\square$ & $\square$ & $\square$ & $\square$ & $\square$ \\
	\midrule
	\multicolumn{9}{@{}l}{\textbf{Teil 2: Umgang mit dem Portal}} \\
	PEOU1 & Es würde mir leichtfallen, den Umgang mit dem Portal zu erlernen. \newline {\scriptsize\itshape Learning to operate [X] would be easy for me.} & $\square$ & $\square$ & $\square$ & $\square$ & $\square$ & $\square$ & $\square$ \\
	PEOU2 & Ich würde es leicht finden, das Portal dazu zu bringen, das zu tun, was ich möchte. \newline {\scriptsize\itshape I would find it easy to get [X] to do what I want it to do.} & $\square$ & $\square$ & $\square$ & $\square$ & $\square$ & $\square$ & $\square$ \\
	PEOU3 & Meine Interaktion mit dem Portal wäre klar und verständlich. \newline {\scriptsize\itshape My interaction with [X] would be clear and understandable.} & $\square$ & $\square$ & $\square$ & $\square$ & $\square$ & $\square$ & $\square$ \\
	PEOU4 & Ich würde den Umgang mit dem Portal als flexibel empfinden. \newline {\scriptsize\itshape I would find [X] to be flexible to interact with.} & $\square$ & $\square$ & $\square$ & $\square$ & $\square$ & $\square$ & $\square$ \\
	PEOU5 & Es würde mir leichtfallen, im Umgang mit dem Portal sicher zu werden. \newline {\scriptsize\itshape It would be easy for me to become skillful at using [X].} & $\square$ & $\square$ & $\square$ & $\square$ & $\square$ & $\square$ & $\square$ \\
	PEOU6 & Ich würde das Portal einfach zu bedienen finden. \newline {\scriptsize\itshape I would find [X] easy to use.} & $\square$ & $\square$ & $\square$ & $\square$ & $\square$ & $\square$ & $\square$ \\
	\midrule
	\multicolumn{9}{@{}l}{\textbf{Teil 3: Künftige Nutzung}} \\
	BI1 & Angenommen, das Portal stünde mir zur Verfügung: Ich sage voraus, dass ich es künftig regelmäßig nutzen werde. & $\square$ & $\square$ & $\square$ & $\square$ & $\square$ & $\square$ & $\square$ \\
	BI2 & Angenommen, das Portal stünde mir zur Verfügung: Ich sage voraus, dass ich es künftig regelmäßig nutzen werde. (Endpunkte vertauscht gedruckt, Wert = 8 minus Rohwert) & $\square$ & $\square$ & $\square$ & $\square$ & $\square$ & $\square$ & $\square$ \\
\end{longtable}
\endgroup
